% TRACE Research Poster — CS 5542 Research-A-Thon 2026
% HOW TO USE: Open the Lethbridge template on Overleaf ("Open as Template"),
%             then replace main.tex with this file. Compile with LuaLaTeX.
% Template: https://www.overleaf.com/latex/templates/university-of-lethbridge-unofficial-poster-template/nddfzgvqvfwf

\documentclass[final]{beamer}

% ====================
% Packages
% ====================

\usepackage[T1]{fontenc}
\usepackage[utf8]{luainputenc}
\usepackage{lmodern}
\usepackage[size=custom, width=91.44, height=121.92, scale=1.15]{beamerposter}
\usetheme{gemini}
\usecolortheme{msu}
\usepackage{graphicx}
\usepackage{booktabs}
\usepackage{tikz}
\usetikzlibrary{arrows.meta, positioning, shapes.geometric, calc}
\usepackage{pgfplots}
\pgfplotsset{compat=1.14}
\usepackage{anyfontsize}
\usepackage{colortbl}
\usepackage{xcolor}

% ====================
% Custom colors
% ====================
\definecolor{tracegreen}{HTML}{2E7D32}
\definecolor{tracered}{HTML}{C62828}
\definecolor{wingreen}{HTML}{C8E6C9}
\definecolor{losered}{HTML}{FFCDD2}
\definecolor{lightblue}{HTML}{BBDEFB}
\definecolor{lightorange}{HTML}{FFE0B2}
\definecolor{lightpurple}{HTML}{E1BEE7}

% ====================
% Lengths
% ====================

\newlength{\sepwidth}
\newlength{\colwidth}
\setlength{\sepwidth}{0.025\paperwidth}
\setlength{\colwidth}{0.3\paperwidth}

\newcommand{\separatorcolumn}{\begin{column}{\sepwidth}\end{column}}

% ====================
% Title
% ====================

\title{TRACE: Terminal Recognition and Attribution via Command Entropy}
\author{Murali Ediga \quad \textit{Advisor: Dr.\ Sudipta Chattopadhyay}}
\institute[shortinst]{ASSET Lab, University of Missouri--Kansas City}

% ====================
% Footer
% ====================

\footercontent{
  CS 5542 Big Data Analytics \& Management \hfill
  Research-A-Thon 2026 \hfill
  \texttt{github.com/muralikrish9/CS5542}}

\begin{document}

\begin{frame}[t]
\begin{columns}[t]
\separatorcolumn

% ═══════════════════════════════════════════════════════════════════════════════
%  COLUMN 1: Problem, Related Work, Dataset, Methodology
% ═══════════════════════════════════════════════════════════════════════════════
\begin{column}{\colwidth}

  \begin{alertblock}{Problem Statement / Motivation}

    \textbf{AI agents can now autonomously hack into systems} --- scanning
    networks, exploiting vulnerabilities, and stealing data with
    \textbf{zero human involvement}. But here's the catch:

    \begin{center}
      \textit{\large ``Every LLM leaves a fingerprint in the commands it types.''}
    \end{center}

    \begin{itemize}
      \item \textbf{7 frontier models} (GPT-5.4, Claude, Gemini, DeepSeek\ldots)
            each have distinct, exploitable command patterns
      \item \textbf{No existing work} identifies LLMs from terminal behavior
            in agentic settings --- LLMmap uses chat, CHeaT uses static text
      \item \textbf{Each model has different weaknesses} --- but blind defenders
            waste their best weapons on the wrong target
    \end{itemize}

    \heading{Goal}

    \textbf{Identify the attacking model from shell commands alone,
    then hack it back} with a targeted prompt injection.

  \end{alertblock}

  \begin{block}{Related Work --- Why Existing Approaches Fall Short}

    \begin{table}
      \centering
      \begin{tabular}{l c c c}
        \toprule
        & \textbf{LLMmap} & \textbf{CHeaT} & \textbf{TRACE (ours)} \\
        \midrule
        Signal source & Chat responses & Generated text & \cellcolor{wingreen}\textbf{Terminal cmds} \\
        Context & Conversational & Static text & \cellcolor{wingreen}\textbf{Agentic tool-use} \\
        Task & Which LLM? & AI vs human & \cellcolor{wingreen}\textbf{7-class family} \\
        Cross-scaffold & N/A & N/A & \cellcolor{wingreen}\textbf{Yes (LOSO)} \\
        Evasion testing & No & No & \cellcolor{wingreen}\textbf{3 strategies} \\
        DPI integration & No & No & \cellcolor{wingreen}\textbf{Targeted routing} \\
        \bottomrule
      \end{tabular}
    \end{table}

    \textbf{LLMmap} (Pasquini et al., 2024): fingerprints LLMs via
    conversational prompts --- but agents execute commands, not chat.
    \textbf{CHeaT} (Wang et al., 2024): binary AI-vs-human detection on static
    text --- not multi-class, not agentic. Neither evaluates cross-scaffold
    transfer or links fingerprinting to active defense.

  \end{block}

  \begin{block}{Dataset Description}

    \heading{Model selection criteria}

    Top-performing on agentic benchmarks (SWE-bench, HumanEval, GAIA),
    capable of complex multi-step tool calls, widely deployed in production.

    \begin{table}
      \centering
      \begin{tabular}{l l r}
        \toprule
        \textbf{Family} & \textbf{Provider} & \textbf{Sessions} \\
        \midrule
        Claude Opus 4.6 & Anthropic  & 249 \\
        GPT-5.4         & OpenAI     & 238 \\
        Gemini 3.1      & Google     & 238 \\
        DeepSeek-V3     & DeepSeek   & 218 \\
        Qwen-2.5        & Alibaba    & 198 \\
        Kimi-K2         & Moonshot   & 198 \\
        GLM-5           & Zhipu AI   & 336 \\
        \midrule
        \multicolumn{2}{l}{\textbf{Grand total}} & \textbf{1,875} \\
        \bottomrule
      \end{tabular}
    \end{table}

    \textbf{3 scaffolds:} Claude Code (CC), PentestGPT (PGPT), ReAct agent.
    Each session = 10--50 commands on a Docker-isolated Linux honeypot
    with standardized CTF challenges.

  \end{block}

  \begin{block}{Methodology / Approach}

    \heading{What failed (v1): hand-crafted behavioral features}

    80 features (deliberation depth, error recovery, command entropy).
    \textbf{Cross-scaffold F1 = 0.11} --- features predicted scaffold, not model.

    \heading{What works (v2): TF-IDF on raw commands}

    \begin{table}
      \centering
      \begin{tabular}{l l}
        \toprule
        \textbf{Stage} & \textbf{Details} \\
        \midrule
        Input & All commands $\rightarrow$ single text string \\
        Vectorizer & TF-IDF: unigrams + bigrams, 50K features \\
        Classifier & LinearSVC ($C=1.0$), one-vs-rest \\
        Evaluation & 5-fold CV, LOSO, evasion, blind test \\
        \bottomrule
      \end{tabular}
    \end{table}

    Each LLM has \textbf{model-intrinsic} command preferences from pretraining
    --- Claude prefers \texttt{python3 -c}, DeepSeek chains with \texttt{\&\&},
    GLM-5 annotates with \texttt{echo}. These transfer across scaffolds.

  \end{block}

\end{column}

\separatorcolumn

% ═══════════════════════════════════════════════════════════════════════════════
%  COLUMN 2: System Architecture, Results
% ═══════════════════════════════════════════════════════════════════════════════
\begin{column}{\colwidth}

  \begin{block}{System Architecture --- Live Deployment Pipeline}

    \begin{center}
    \begin{tikzpicture}[
        node distance=0.5cm,
        >={Stealth[length=3.5mm, width=2mm]},
        box/.style={rectangle, rounded corners=3pt, draw=black!50, fill=#1,
                    line width=0.8pt, minimum width=12cm, minimum height=1.2cm,
                    align=center, font=\small},
        arrow/.style={->, line width=1pt, color=black!50},
        lbl/.style={font=\scriptsize\itshape, text=black!60},
      ]
      \node[box=lightpurple] (agent)
        {\textbf{AI Agent Session} --- unknown model family};
      \node[box=lightorange, below=of agent] (hp)
        {\textbf{Honeypot Server} --- captures raw commands};
      \node[box=lightblue, below=of hp] (agg)
        {\textbf{Session Aggregator} --- groups commands per session};
      \node[box=wingreen, below=of agg] (clf)
        {\textbf{TF-IDF + LinearSVC} --- predict family + confidence};
      \node[box=losered, below=of clf] (dpi)
        {\textbf{Targeted DPI Router} --- deploy best payload};

      \draw[arrow] (agent) -- (hp) node[midway, right, lbl] {SSH session};
      \draw[arrow] (hp) -- (agg) node[midway, right, lbl] {command stream};
      \draw[arrow] (agg) -- (clf) node[midway, right, lbl] {session text};
      \draw[arrow] (clf) -- (dpi) node[midway, right, lbl] {e.g.\ ``deepseek 94\%''};
    \end{tikzpicture}
    \end{center}

    Deployed on GCP. Live demo: Opus identified at \textbf{87\%} (DPI-resistant),
    DeepSeek at \textbf{94\%} $\rightarrow$ routed to Vanilla payload.

  \end{block}

  \begin{block}{Results \& Evaluation}

    \heading{Fingerprinting accuracy (5-fold CV)}

    \begin{center}
      {\fontsize{42}{48}\selectfont \textbf{F1 = 0.979}} \quad {\large $\pm$ 0.003}
    \end{center}

    \heading{Scaffold generalization (LOSO)}

    \begin{table}
      \centering
      \begin{tabular}{l c l}
        \toprule
        \textbf{Held-out} & \textbf{Macro-F1} & \\
        \midrule
        CC (Claude Code)  & 0.783 & Moderate transfer \\
        PGPT (PentestGPT) & \textbf{0.953} & Strong transfer \\
        ReAct              & 0.602 & Hardest (unique framing) \\
        \midrule
        \textbf{Mean} & \textbf{0.812} & \textit{7.4$\times$ over v1 features} \\
        \bottomrule
      \end{tabular}
    \end{table}

    \heading{Classifier comparison}

    \begin{table}
      \centering
      \begin{tabular}{l c c}
        \toprule
        \textbf{Classifier} & \textbf{CV F1} & \textbf{LOSO F1} \\
        \midrule
        \rowcolor{wingreen} \textbf{LinearSVC (ours)} & \textbf{0.979} & \textbf{0.812} \\
        Logistic Regression & 0.975 & 0.799 \\
        SGD (hinge) & 0.971 & 0.788 \\
        Random Forest & 0.941 & 0.682 \\
        Multinomial NB & 0.935 & 0.710 \\
        \bottomrule
      \end{tabular}
    \end{table}

    LinearSVC wins on both in-distribution \emph{and} cross-scaffold transfer.

    \heading{Cross-model generalization}

    \begin{table}
      \centering
      \begin{tabular}{l l c}
        \toprule
        \textbf{Variant} & \textbf{Target family} & \textbf{Accuracy} \\
        \midrule
        Opus 4.6 (control) & claude\_opus & \cellcolor{wingreen}100\% \\
        Sonnet 4.6 $\rightarrow$ Opus & claude\_opus & \cellcolor{lightblue}\textbf{72\%} \\
        Gemini 2.5 $\rightarrow$ 3.1 & gemini31 & \cellcolor{lightblue}\textbf{68\%} \\
        \bottomrule
      \end{tabular}
    \end{table}

    \heading{Feature ablation}

    \begin{table}
      \centering
      \begin{tabular}{l c c}
        \toprule
        \textbf{Condition} & \textbf{CV F1} & \textbf{LOSO} \\
        \midrule
        \textbf{Full bigram (ours)} & \textbf{0.979} & \textbf{0.812} \\
        Verb-only unigram & 0.829 & 0.568 \\
        \rowcolor{losered} Random (1/7) & 0.143 & 0.143 \\
        \bottomrule
      \end{tabular}
    \end{table}

  \end{block}

\end{column}

\separatorcolumn

% ═══════════════════════════════════════════════════════════════════════════════
%  COLUMN 3: DPI Application, Demo, Conclusion
% ═══════════════════════════════════════════════════════════════════════════════
\begin{column}{\colwidth}

  \begin{alertblock}{Hacking the Hacker: Targeted Defensive Prompt Injection}

    \textbf{If an AI agent is hacking you, hack it back.}
    TRACE identifies the attacking model, then exploits \emph{its}
    specific weaknesses with a tailored prompt injection.

    Three DPI payloads planted in environment files the agent reads:

    \begin{table}
      \centering
      \footnotesize
      \begin{tabular}{l l}
        \toprule
        \textbf{Payload} & \textbf{Mechanism} \\
        \midrule
        Vanilla & Heredoc in \texttt{WELCOME.txt} writes task briefing \\
        M2 Authority & Impersonates forensic audit demand \\
        FC Coercion & Frames log entry as format violation \\
        \bottomrule
      \end{tabular}
    \end{table}

    \heading{ASR (\%) --- every family has a different weakness}

    \begin{table}
      \centering
      \begin{tabular}{l c c c c}
        \toprule
        \textbf{Family} & \textbf{Vanilla} & \textbf{M2} & \textbf{FC} & \textbf{Best} \\
        \midrule
        \rowcolor{wingreen} Claude & 0\% & 0\% & 0\% & \textbf{0\%} \\
        DeepSeek & \cellcolor{losered}\textbf{95\%} & 87\% & 75\% & 95\% \\
        Gemini & \cellcolor{losered}85\% & 70\% & 60\% & 85\% \\
        GPT-5.4 & 13\% & 17\% & \cellcolor{losered}\textbf{90\%} & 90\% \\
        GLM-5 & 53\% & 80\% & \cellcolor{losered}\textbf{90\%} & 90\% \\
        Kimi & 41\% & 57\% & \cellcolor{losered}\textbf{88\%} & 88\% \\
        Qwen & 73\% & 62\% & 57\% & 73\% \\
        \bottomrule
      \end{tabular}
    \end{table}

    \heading{Blind vs.\ targeted defense}

    \begin{table}
      \centering
      \begin{tabular}{l c}
        \toprule
        \textbf{Strategy} & \textbf{Mean ASR} \\
        \midrule
        Blind (best single payload) & 60.3\% \\
        \rowcolor{wingreen} \textbf{Fingerprint $\rightarrow$ route to best} & \textbf{86.8\%} \\
        \midrule
        \textbf{Targeting lift} & \textbf{+26.5pp} \\
        \bottomrule
      \end{tabular}
    \end{table}

    \textbf{Know thy enemy} $\rightarrow$ choose the right weapon $\rightarrow$
    \textbf{+26.5pp lift} over blind defense.

  \end{alertblock}

  \begin{block}{Demo / Live Deployment}

    TRACE runs as a live service on a GCP VM:
    \begin{itemize}
      \item \textbf{Cowrie honeypot} captures incoming SSH agent sessions
      \item \textbf{Real-time classifier} aggregates commands in 2-min windows
      \item \textbf{Dashboard} displays predicted family + confidence
      \item \textbf{DPI router} selects and deploys the optimal payload
    \end{itemize}

    Live test: an Opus-driven Claude Code agent was correctly identified at
    \textbf{87\% confidence} and flagged as DPI-resistant. A DeepSeek ReAct
    agent was identified at \textbf{94\%} and routed to the Vanilla payload.

  \end{block}

  \begin{alertblock}{Conclusion \& Future Work}

    \heading{Key contributions}

    \begin{enumerate}
      \item \textbf{First terminal-level LLM fingerprinting} --- fills the gap
            left by LLMmap (chat) and CHeaT (text)
      \item \textbf{98\% accuracy} across 7 flagship models, 3 scaffolds
      \item \textbf{Scaffold-invariant} (LOSO F1 = 0.812) --- generalizes to unseen frameworks
      \item \textbf{Hack the hacker: +26.5pp targeting lift} via fingerprint-guided DPI
      \item \textbf{Deployed live} on GCP honeypot
    \end{enumerate}

    \heading{Future work}

    \begin{itemize}
      \item \textbf{Intelligence layer:} fine-tuned Qwen for low-confidence
            session resolution
      \item \textbf{Multi-agent swarms:} disambiguate concurrent models
      \item \textbf{Temporal features:} inter-command timing analysis
    \end{itemize}

  \end{alertblock}

  \begin{block}{References}

    \footnotesize
    \begin{enumerate}
      \item Pasquini et al. ``LLMmap: Fingerprinting For LLMs.'' USENIX Security, 2024.
      \item Wang et al. ``CHeaT: A Chatbot Detection Tool.'' USENIX Security, 2024.
      \item Fang et al. ``LLM Agents Can Autonomously Hack Websites.'' arXiv:2402.06664, 2024.
      \item Greshake et al. ``Indirect Prompt Injection.'' AISec, 2023.
      \item Deng et al. ``PentestGPT: An LLM-empowered Automatic Penetration Testing Tool.'' USENIX Security, 2024.
    \end{enumerate}

  \end{block}

\end{column}

\separatorcolumn
\end{columns}
\end{frame}

\end{document}
